%=======================================================================================
% METHODOLOGY CHAPTER -- 2026-08-13, 12:30
% Restructured from methodology_merged_2026-08-11_2221.tex per the structural review notes of
% 2026-08-13. Content (equations, numbers, model list) is unchanged from that version except
% where a fix below specifically required a correction; this pass only reorganises the
% hierarchy, renames headings to match what each part actually contains, and removes code-level
% identifiers (file names, script/class/column names) in favour of plain descriptions. No
% performance results, interpreted associations, or model-comparison verdicts appear here --
% see the Results chapter. Sentences are left terse/skeletal in places deliberately -- fill in
% final prose by hand. Points still flagged "% TODO" were flagged in the prior version and have
% not yet been resolved.
%=======================================================================================

\chapter{Methodology and Experimental Design}
\label{ch:methodology}

This chapter describes how the raw plot data (Chapter~\ref{ch:data}) were converted into the
results used to answer the project's three research questions. It is organised by method, not by
the order experiments were run: the experimental framework first (\S\ref{sec:framework}), then
common design decisions and leakage control (\S\ref{sec:common-design}), environmental
feature-set construction (\S\ref{sec:feature-sets}), the top-height prediction models
(\S\ref{sec:rq1}), the two attribution analyses (\S\ref{sec:attribution}), evaluation design
(\S\ref{sec:evaluation}), and finally assumptions/limitations and reproducibility
(\S\ref{sec:limitations}--\S\ref{sec:reproducibility}).

\section{Research questions and experimental framework}
\label{sec:framework}

% Short paragraph: prediction (RQ1) vs attribution (RQ2, RQ3) as the two analytical modes of the
% chapter; RQ2 and RQ3 share a starting point (the pooled Chapman-Richards curve) but diverge in
% what they attribute -- pooled residual vs plot-specific asymptote.

RQ1 is not a single model competition. It tests, together: predictive accuracy of the
top-height models; whether environmental/terrain input improves prediction over no-environment
input; whether Chapman--Richards guidance improves prediction over an unguided network of the
same architecture; and whether any of those conclusions hold under spatial and temporal
generalisation, not only under a random/plot-level split.

\begin{table}[!htbp]
\centering
\scriptsize
\renewcommand{\arraystretch}{1.2}
\begin{tabular}{p{0.8cm}p{1.4cm}p{3.3cm}p{4.2cm}p{2.9cm}}
\toprule
RQ & Unit & Target & Model families compared & Evaluation design \\
\midrule
RQ1 & Plot-survey row & 95th-percentile LiDAR top height & Chapman--Richards curve, average-by-age
lookup, linear/random forest/XGBoost, deep neural network, Chapman--Richards-guided neural network
variants & Spatial-block (primary), pooled spatial $k$-fold, temporal, plot-level \\
RQ2 & Plot & Mean residual from the single pooled Chapman--Richards curve, across a plot's survey
history & Elastic Net, XGBoost, two-stage mixed-effects model & Spatial-block \\
RQ3 & Plot & Difference between a plot's own fitted asymptote and its yield-class-implied
asymptote & Elastic Net, XGBoost, GNNWR & Spatial-block \\
\bottomrule
\end{tabular}
\caption{Research questions mapped to unit of analysis, target, model families, and evaluation
design.}
\label{tab:rq-framework}
\end{table}

\section{Common experimental design and leakage control}
\label{sec:common-design}

\subsection{Cohort roles and units of analysis}
\label{sec:cohorts}

Two overlapping cohorts are used throughout: 4survey (four repeated LiDAR surveys per plot) and
6survey (six repeated surveys, a strict subset of 4survey's plots with a longer, denser survey
history but a smaller training set). 4survey is the primary cohort for every research question;
6survey is run in parallel as a secondary check wherever the sample size permits it. RQ1's unit of
analysis is the plot-survey row; RQ2 and RQ3 aggregate to one row per plot (Table~\ref{tab:rq-framework}).
Dataset-level detail (plot counts, survey years, coverage) is given in Chapter~\ref{ch:data} and
is not repeated here.

\subsection{Spatial and temporal evaluation}
\label{sec:splits}

Four split types, each answering a different generalisation question:

\begin{itemize}
\setlength\itemsep{0pt}
\item \textbf{Spatial block} (primary). Whole forestry compartments are assigned to
  train/validation/test so that no compartment appears on both sides of a split -- a plot-level or
  random split would let a model see a near-identical neighbouring plot during training,
  overstating how well it generalises to genuinely new ground. Compartments are chosen by
  row-count-balanced random assignment (60/20/20 train/validation/test, seed 42), with a 60\,m
  buffer around every validation/test plot: any training plot within 60\,m of a held-out plot is
  excluded from training (validation/test plots are never removed). The buffer distance is set
  from typical plot spacing rather than from a measured spatial-autocorrelation range; leakage
  protection at a larger spatial scale comes from holding out whole compartments, not from the
  buffer alone \citep{roberts2017}. % TODO: citation -- Roberts et al. 2017; no .bib entry yet.
  % TODO: the shared-curve residual's own semivariogram range (~3,956 m) is an observed
  % diagnostic result -- confirm whether it belongs here (if estimated on training data only and
  % used to justify the buffer) or should move to Results/Appendix.
\item \textbf{Spatial block, pooled $k$-fold} (stronger confirmation of the single spatial-block
  split). Compartments are assigned to 5 row-count-balanced folds; each fold in turn is held out
  as test, the next as validation, the remainder as train, so every compartment is evaluated
  exactly once. Predictions from all 5 folds are pooled into one population-wide evaluation set
  rather than averaged as 5 separate summary statistics.
\item \textbf{Temporal} (secondary, harder test). Whole survey years, not compartments, are
  assigned to train/validation/test: 4survey trains on 2008 and 2012, validates on 2021, tests on
  2023; 6survey trains on 2002/2006/2008/2012, validates on 2021, tests on 2023. This asks whether
  a model trained on earlier growth still predicts a real future survey it never saw -- a
  different, harder question than spatial generalisation. A narrow-gap variant (training years
  closer to the test year) is also defined as a secondary check, with validation deliberately
  assigned to the earliest available year rather than a middle year.
  % TODO: the reason for that choice (a failed smoke test with a middle validation year) is
  % implementation history -- move to Appendix or omit from the main text.
\item \textbf{Plot level} (sanity floor only, not a generalisation claim). Individual plots are
  split at random (60/20/20, seed 42) with no spatial holdout at all -- kept only as an easy-case
  reference, since a held-out plot's nearest neighbour is typically a training plot metres away.
\end{itemize}

\subsection{Training, validation and test separation}
\label{sec:train-val-test}

Every model in this project reads the same split assignment for a given (cohort, split type,
seed) combination, so results are always compared on identical train/validation/test membership.
Compartment/block/sub-compartment identifier columns are never included in any model's feature
list, enforced automatically. Shared data partitions and matched architectures reduce confounding
when comparing model families under the same split -- e.g. the DNN and the Chapman--Richards-guided
network (\S\ref{sec:rq1-dnn}--\S\ref{sec:rq1-pinn}) -- though residual differences can still
reflect optimisation behaviour or stochastic variation rather than the model design alone.

\subsection{Preprocessing and leakage prevention}
\label{sec:preprocessing}

Continuous features are standardised (zero mean, unit variance) using a scaler fitted on the
training partition only, then applied unchanged to validation/test, so no information about
held-out rows leaks into the scaling itself. Age is scaled with its own separate scaler, kept
apart from other continuous features so its training-split standard deviation can be read back out
directly for the physics loss's unit correction (\S\ref{sec:rq1-pinn-loss}). Unordered categorical
variables (three soil-class rasters) are one-hot encoded, never treated as an ordered numeric
scale. Rows with a missing value in any requested feature are dropped whole-plot (every survey
year of that plot, not just the affected row), so no model ever sees a partial trajectory for a
plot it has otherwise seen.

\section{Environmental feature-set construction}
\label{sec:feature-sets}

Three questions each need their own screened list of candidate environmental/management
variables: RQ1's predictive tiers, RQ2's shared-curve attribution set, and RQ3's plot-level
attribution set. All three follow the same recipe below, differing only in candidate pool and
target column. Full per-variable diagnostics, VIF tables, and complete final column lists are
given in Appendix~\ref{app:feature-sets}; this section states the rules only.

\subsection{Candidate variables and protected baseline features}
\label{sec:feature-baseline}

Four stand/management columns (canopy cover, time since thinning, a paired missingness flag, and
recent-thinning fraction) are present unconditionally in every set and protected from removal at
every later screening stage. A fifth candidate, thinning fraction, was confirmed to be an exact
linear complement of the missingness flag and is not included separately.

\subsection{Deduplication, ranking and multicollinearity screening}
\label{sec:feature-screening}

\begin{enumerate}
\setlength\itemsep{0pt}
\item \textbf{Deduplication.} Deterministic duplicates (a variable that is a documented closed-form
  function of another candidate already in the pool) are dropped outright; near-exact empirical
  duplicates (Spearman $|\rho|\geq0.95$) are reduced to one representative per pair, preferring an
  externally measured variable over a derived one.
\item \textbf{Importance ranking.} Three signals are computed against each avenue's own target and
  combined by rank-averaging: univariate Spearman $|\rho|$; permutation importance; and
  drop-column ablation (the held-out $R^2$ change from removing one variable, given every other
  candidate already in the model). Each signal has a known blind spot -- correlation ignores
  context, permutation is biased under correlated predictors, ablation depends on one model/split
  -- so this is a soft preference for cross-signal support: a variable that scores strongly on one
  signal but weakly on the others can still reach a middling combined rank, not a hard
  multi-signal gate.
\item \textbf{Variance-inflation screening.} Iterative VIF reduction (threshold 5.0
  \citep{dormann2013}) is applied on top of ranking, with the protected baseline columns exempt
  from removal (they remain in the design matrix, affecting every other column's VIF, but are
  never themselves dropped).
\end{enumerate}

% TODO: state explicitly which partition (training only, vs. the full pre-split pool) the
% ranking/VIF screening was computed on; whether ranking was repeated separately inside each
% spatial fold or computed once; and whether Set1-Set5 were frozen before the final evaluation
% split was drawn, so a reader can confirm the test compartments had no influence on feature-set
% construction.

\subsection{Research-question-specific feature sets}
\label{sec:feature-final}

Set1 = baseline only. Set2 = baseline + the top 10 candidates by combined rank, itself
VIF-screened with backfill from the next-ranked candidate. Set3 = baseline + terrain/wind
candidates in the upper half of their own category by combined rank, VIF-screened. Set4 = Set3
plus every other category's upper-half candidates, VIF-screened. Set5 = baseline plus every
deduplicated candidate, VIF-screened, no rank filter.

\begin{table}[!htbp]
\centering
\scriptsize
\renewcommand{\arraystretch}{1.2}
\begin{tabular}{p{2.3cm}p{4.2cm}p{1.6cm}p{2cm}p{3.4cm}}
\toprule
Feature set family & Feeds & Model comparison & Attribution & Purpose \\
\midrule
RQ1 Set1--Set5 & DNN and Chapman--Richards-guided network variants (environment-conditioned) &
Primary & No & DNN-vs-guided-network comparison under progressively richer terrain/wind/climate
input; no coefficient-based model reads these lists. \\
RQ2 Set1--Set5 & Elastic Net, XGBoost, two-stage mixed-effects model & Secondary & Primary & Which
conditions align with persistent departure from the shared curve. \\
RQ3 Set1--Set5 & Elastic Net, XGBoost, GNNWR & Secondary & Primary & Which conditions align with
plot-specific trajectory deviation; GNNWR additionally tests whether that association is
spatially varying. \\
\bottomrule
\end{tabular}
\caption{Purpose of the final feature sets by research question. Full column lists: Appendix~\ref{app:feature-sets}.}
\label{tab:feature-set-purpose}
\end{table}

% FIGURE PLACEHOLDER (pipeline overview -- simplified, high level only)
% Caption: "From plots to results: cohorts feed spatial/temporal partitions and a shared
% screening process producing Set1-5 per research question. Three branches follow: RQ1's
% predictive models, RQ2's shared-curve residual attribution, RQ3's plot-specific asymptote
% attribution, each with its own evaluation output."
%
% Confirmed evaluation output per branch (do not claim more than this without checking the
% notebooks under notebooks/spatial_analysis/, which were not inspected for this pass): RQ1
% produces per-fold and pooled predictive metrics under each split type; RQ2 produces a
% fixed-effects coefficient table plus held-out residual predictions; RQ3 produces per-fold
% predictive metrics, a per-plot coefficient table (GNNWR), and a residual spatial-autocorrelation
% diagnostic. No rendered coefficient/residual map was confirmed anywhere in the modelling
% pipeline -- do not draw one as an output box unless separately verified.
%
% Mermaid sketch (simplified, for redrawing):
% ```mermaid
% flowchart TD
%     A["Cohorts: 4survey / 6survey"] --> B["Spatial and temporal<br/>partitions"]
%     A --> C["Candidate environmental<br/>and management variables"]
%     C --> D["RQ-specific feature<br/>screening: Set1-5"]
%     B --> E1
%     B --> E2
%     B --> E3
%     D --> E1["RQ1: top-height<br/>prediction models"]
%     D --> E2["RQ2: shared-curve<br/>residual attribution"]
%     D --> E3["RQ3: plot-specific<br/>asymptote attribution"]
%     E1 --> F1["Predictive metrics,<br/>per-fold and pooled"]
%     E2 --> F2["Fixed-effects coefficient<br/>table + held-out residuals"]
%     E3 --> F3["Predictive metrics +<br/>per-plot coefficient table"]
% ```
\begin{figure}[!htbp]
\centering
\fbox{\parbox{0.9\textwidth}{\centering\vspace{4cm}FIGURE PLACEHOLDER --- pipeline overview, see
comment above source for exact content\vspace{4cm}}}
\caption{Overview of the full pipeline from raw data to the three research questions' results.}
\label{fig:pipeline-overview}
\end{figure}

\section{Top-height prediction models}
\label{sec:rq1}

Target: the 95th-percentile LiDAR-derived top height, per plot per survey. Six model families are
compared, sharing splits, metrics, and (where applicable) network architecture.

\subsection{Reference and conventional comparison models}
\label{sec:rq1-baselines}

\textbf{Chapman--Richards (CR) curve.} A three-parameter curve
\begin{equation}
H(a) = y_{\max}\left(1-e^{-ka}\right)^{p}
\label{eq:cr}
\end{equation}
is fitted by nonlinear least squares (five restarts, lowest sum-of-squared-error kept), where
$H(a)$ is predicted top height at age $a$ (years), $y_{\max}$ is the asymptotic maximum height,
$k$ controls growth rate, and $p$ is a shape parameter. This is a single \emph{pooled},
population-average curve -- one fit per cohort, not per plot or per compartment -- fitted on the
training partition of the split in use (never on validation/test rows), refitted separately per
fold under the pooled spatial-block $k$-fold split so that a given fold's held-out compartments
never influence its own frozen curve. The pooled curve's fitted $(y_{\max},k,p)$ are treated as
constants downstream (\S\ref{sec:rq1-pinn-loss}), never as trainable parameters.

\textbf{Average-by-age.} A lookup table of mean observed height per integer age year, computed on
the same training rows as the CR fit; an unseen age falls back to the overall training mean
height.

\textbf{Linear, random forest, and XGBoost baselines.} Conventional supervised baselines fit
directly on the no-environment feature set, evaluated under the same splits as every other model.
XGBoost's hyperparameters (max depth $\in\{2,3,4,6\}$, $\ell_2$ regularisation
$\in\{1,5,20\}$, 500 boosting rounds fixed with early stopping, learning rate 0.1 fixed) are
selected by validation-set $R^2$ over this 12-combination grid; full grid and selection results:
Appendix~\ref{app:hyperparameters}.
% TODO: verify -- confirm whether the environmental-feature-set XGBoost fits reported in RQ1's
% sweep used this same grid/selection procedure or a separate configuration.

\subsection{Deep neural network}
\label{sec:rq1-dnn}

A single shared architecture is used for every DNN and Chapman--Richards-guided variant in this
project: three hidden layers of 128 units each, leaky-ReLU activation, age concatenated with every
other input feature. Reusing one architecture throughout means any difference between the DNN and
the guided network's results is attributable to the loss function, not to a different network. The
DNN's loss is plain mean squared error,
\begin{equation}
\mathcal{L}_{\text{data}} = \frac{1}{N}\sum_{i=1}^{N}\left(\hat H_i - H_i\right)^2,
\label{eq:data-loss}
\end{equation}
where $\hat H_i$ is the network's predicted height for row $i$ and $H_i$ is the observed height,
plus an L1 weight penalty ($\lambda_{L1}=10^{-5}$). The environment-conditioned DNN is
architecturally identical to the no-environment DNN; its terrain/wind block is simply concatenated
into the same input, widening that vector rather than adding a second network.

\subsection{Chapman--Richards-guided neural network}
\label{sec:rq1-pinn}

Three physics-informed variants share the DNN's exact network class, differing only in loss terms
and, for the environment-conditioned variants, in an added small sub-network. The no-environment
variant uses the CR curve's $(y_{\max},k,p)$ as the single pooled, frozen cohort constants from
Equation~\ref{eq:cr}.

\subsubsection{Physics and trajectory losses}
\label{sec:rq1-pinn-loss}

Every guided variant optimises the same three-term objective, using $H'(a)$, the CR curve's
analytical derivative,
\begin{equation}
H'(a) = y_{\max}\,p\,k\,e^{-ka}\left(1-e^{-ka}\right)^{p-1}.
\label{eq:cr-derivative}
\end{equation}
The \emph{physics} term matches the network's own instantaneous growth rate, obtained by automatic
differentiation of its prediction with respect to age, against the analytical rate at that same
age:
\begin{equation}
\mathcal{L}_{\text{phys}} = \frac{1}{N}\sum_{i=1}^{N}\left(\frac{\partial \hat H_i}{\partial a_i} -
H'(a_i)\right)^2.
\label{eq:physics-loss}
\end{equation}
Because age is standardised for network input (\S\ref{sec:preprocessing}), the automatic-differentiation
derivative is with respect to standardised age; it is rescaled by the training-split age standard
deviation before comparison against $H'(a_i)$, which is defined in real age units.

The \emph{trajectory} term is a separate, finite-difference check: for $M$ pairs of real,
consecutive same-plot survey observations (both endpoints drawn only from that plot's own
training-labelled rows), the network's own predicted height \emph{change} between the pair must
match the analytical rate at the pair's mid-age,
\begin{equation}
\mathcal{L}_{\text{traj}} = \frac{1}{M}\sum_{j=1}^{M}\left(\frac{\hat H_{j,\text{later}} -
\hat H_{j,\text{earlier}}}{\Delta a_j} - H'(\bar a_j)\right)^2,
\label{eq:trajectory-loss}
\end{equation}
where $\Delta a_j$ is the elapsed time between the pair's two surveys and $\bar a_j$ is their
mid-age. The two terms are deliberately distinct checks: an instantaneous derivative at one point
versus a measured change across two real observations. The full training objective is
\begin{equation}
\mathcal{L} = \mathcal{L}_{\text{data}} + \lambda_{\text{phys}}\mathcal{L}_{\text{phys}} +
\lambda_{\text{traj}}\mathcal{L}_{\text{traj}} + \lambda_{L1}\sum_{\theta}|\theta|,
\label{eq:total-loss}
\end{equation}
with $\lambda_{\text{phys}}=\lambda_{\text{traj}}=1.0$ as the default weighting
(\S\ref{sec:rq1-ablation} for the physics-weight ablation). $(y_{\max},k,p)$ are always plain
constants at every point in this objective, never tensors that accumulate a gradient -- training
updates only the network (and, where present, the sub-networks below), never the process model
itself. A trajectory pair is only formed from two rows both labelled training under the split in
use, so no validation or test survey year can leak into this loss term.

\subsubsection{Environment-conditioned growth parameters}
\label{sec:rq1-pinn-env}

\begin{itemize}
\setlength\itemsep{0pt}
\item The terrain-conditioned variant feeds terrain/wind features into a small sub-network
  (16 hidden units) that outputs a per-plot \emph{additive} adjustment to the pooled $y_{\max}$:
  \begin{equation}
  y_{\max}^{(i)} = y_{\max} + f_{y}\!\left(\mathbf{x}^{\text{terrain}}_i\right),
  \label{eq:ymax-adjust}
  \end{equation}
  where $f_y$ is the sub-network and $\mathbf{x}^{\text{terrain}}_i$ is plot $i$'s terrain/wind
  feature vector. $k$ and $p$ remain the pooled global constants. This sub-network never sees age
  or the no-environment features, and its output is initialised near zero so training starts from
  the known-good pooled curve rather than an arbitrary value. Terrain/wind conditions only the
  growth \emph{ceiling} in this design, not the trajectory shape -- matching the
  asymptote-varies, shape-fixed framing of site-index modelling
  \citep{socha2021}. % TODO: citation -- Socha et al. 2021 (ADA/GADA), no .bib entry yet.
\item A second variant adds a further sub-network producing a per-plot \emph{multiplicative},
  log-space adjustment to $k$:
  \begin{equation}
  k^{(i)} = k \cdot \exp\!\left(f_{k}\!\left(\mathbf{x}^{\text{terrain}}_i\right)\right),
  \label{eq:k-adjust}
  \end{equation}
  again initialised near zero so $k^{(i)}\approx k$ at the start of training. $p$ stays a pooled
  global constant in every variant; no variant makes the curve's shape parameter site-specific.
\end{itemize}

\subsection{Physics-loss controls and ablation}
\label{sec:rq1-ablation}

A physics-weight ablation over $\lambda_{\text{phys}},\lambda_{\text{traj}}\in\{0,1,2\}$ is run
against the default weighting above to test whether Chapman--Richards guidance itself improves
prediction, rather than only whether a model labelled as guided performs well; settings and full
results: Results chapter, not repeated here. The no-environment DNN (\S\ref{sec:rq1-dnn}), sharing
the identical network architecture and trained under the same splits, serves as the matched
control against which the guided network's physics/trajectory terms are isolated -- the two
differ only in loss function, not in network capacity or input (for the no-environment pairing).

% FIGURE PLACEHOLDER (model architecture)
% Caption: "Shared DNN/guided-network architecture. The plain data loss (DNN, and the guided
% network's own data term) uses the main trunk's output directly. The guided network adds two
% consistency terms against the Chapman-Richards curve's analytical derivative: an instantaneous
% check via automatic differentiation (physics loss) and a finite-difference check across real
% survey pairs (trajectory loss). In the environment-conditioned variants, terrain/wind features
% feed only the small asymptote (and, in the rate-conditioned variant, rate) sub-networks --
% never the main trunk -- producing a per-plot adjustment to the CR curve's asymptote/rate that
% the physics and trajectory losses are then computed against."
% Content: a block diagram. Left: [Age, no-env features] -> main trunk (3x128, shared) ->
% [predicted height]. Below/beside: [Terrain/wind features] -> small asymptote sub-network (16
% units) -> [y_max adjustment, additive] and, dashed/optional for the rate-conditioned variant,
% -> rate sub-network -> [k adjustment, multiplicative, log-space]. A separate branch shows
% [predicted height] and [age] feeding into an autograd box labelled dH/da, compared against a
% box labelled "CR analytical derivative (frozen y_max, k, p)" to produce the physics loss; a
% second small diagram shows two real survey observations of one plot feeding two forward
% passes, differenced and compared to the same CR derivative at mid-age, to produce the
% trajectory loss. These survey pairs are pre-built once from real consecutive-survey rows before
% training starts (not synthesised per batch) -- draw them as a small fixed table/lookup feeding
% the diagram, not as an on-the-fly sampling step.
\begin{figure}[!htbp]
\centering
\fbox{\parbox{0.9\textwidth}{\centering\vspace{4cm}FIGURE PLACEHOLDER --- model architecture, see
comment above source for exact content\vspace{4cm}}}
\caption{Shared DNN/Chapman--Richards-guided architecture and the two physics-consistency loss
terms.}
\label{fig:model-architecture}
\end{figure}

\begin{table}[!htbp]
\centering
\scriptsize
\renewcommand{\arraystretch}{1.2}
\begin{tabular}{p{3.2cm}p{3cm}p{2.6cm}p{4.5cm}}
\toprule
Model & Environmental input & Loss terms & Notes \\
\midrule
CR (baseline) & None & N/A (curve fit) & Pooled per cohort, refit per fold for pooled $k$-fold. \\
Average-by-age & None & N/A (lookup) & Same training rows as CR. \\
Linear / random forest / XGBoost & No-environment feature set & Model-specific & XGBoost
hyperparameters selected by validation $R^2$ (Appendix~\ref{app:hyperparameters}). \\
DNN (no environment) & None & Data (Eq.~\ref{eq:data-loss}) + L1 & Shared architecture with every
guided variant. \\
DNN (environment-conditioned) & Terrain/wind, concatenated into main input & Data + L1 & Same
network class, wider input only. \\
Guided network (no environment) & None & Data + physics + trajectory + L1 & $y_{\max},k,p$
pooled, frozen. \\
Guided network (terrain-conditioned asymptote) & Terrain/wind, asymptote sub-network only & Data +
physics + trajectory + L1 & Eq.~\ref{eq:ymax-adjust}; $k,p$ pooled. \\
Guided network (terrain-conditioned asymptote and rate) & Terrain/wind, asymptote and rate
sub-networks & Data + physics + trajectory + L1 & Eqs.~\ref{eq:ymax-adjust}--\ref{eq:k-adjust};
$p$ pooled. \\
\bottomrule
\end{tabular}
\caption{Model family comparison for top-height prediction (RQ1). Full hyperparameters:
Appendix~\ref{app:hyperparameters}.}
\label{tab:rq1-models}
\end{table}

\section{Environmental attribution analyses}
\label{sec:attribution}

Both avenues below start from the single pooled Chapman--Richards curve fitted in
Equation~\ref{eq:cr} but diverge in what they attribute: RQ2 pools every plot's residual against
that one shared curve; RQ3 fits each plot its own asymptote from only that plot's own repeated
observations. Both report association, not causal, relationships.

\subsection{Persistent departure from a shared growth curve (RQ2)}
\label{sec:rq2}

This avenue asks which environmental/management conditions align with a plot's persistent
position above or below the single shared curve, using only the 4survey cohort (6survey's 47
compartments were judged too few for a reliable spatial-block attribution split).

\subsubsection{Target construction}
\label{sec:rq2-target}

For each row $i$ of a plot's own survey history, a residual is computed against the pooled curve,
\begin{equation}
r_i = H_i - H(a_i),
\label{eq:cr-residual}
\end{equation}
and the target used for attribution is the plot's own mean residual across its $T$ observed
survey rows,
\begin{equation}
\overline{r}_{\text{plot}} = \frac{1}{T}\sum_{t=1}^{T} r_{i,t}.
\label{eq:mean-cr-residual}
\end{equation}
Age is excluded from every RQ2 candidate feature: it is the CR curve's only input and therefore
the sole construction variable of $\overline{r}_{\text{plot}}$, making it circular for this
target.
% TODO: verify -- confirm whether any plot contributing to Equation~\ref{eq:mean-cr-residual}'s
% average (computed over all of that plot's survey rows) falls in the test partition of the
% downstream attribution split below; the target is constructed before that split is drawn, which
% is expected rather than a leakage bug, but is not yet explicitly documented as checked.

\subsubsection{Global and mixed-effects attribution models}
\label{sec:rq2-models}

Three models are fitted on $\overline{r}_{\text{plot}}$ using RQ2's Set1--Set5
(\S\ref{sec:feature-final}), under the same spatial-block design as \S\ref{sec:splits}:

\begin{itemize}
\setlength\itemsep{0pt}
\item \textbf{Elastic Net} -- continuous features standardised on the training split; unordered
  soil-class variables one-hot encoded; regularisation strength and $\ell_1$/$\ell_2$ mixing
  ratio selected by 5-fold cross-validation on the training partition only (never touching
  validation or test).
\item \textbf{XGBoost} -- same hyperparameter-selection approach as \S\ref{sec:rq1-baselines}.
\item \textbf{Two-stage mixed-effects model.} Stage 1 is the already frozen, pooled CR curve
  above; Stage 2 is a \emph{linear} mixed model fitted on the residual $\overline{r}_{\text{plot}}$,
  with the candidate environmental variables as fixed effects and a random intercept grouped by
  forestry compartment,
  \begin{equation}
  \overline{r}_{\text{plot}} = \boldsymbol{\beta}^{\top}\mathbf{x}_{\text{plot}} +
  u_{\text{cpmt}} + \varepsilon_{\text{plot}},
  \label{eq:mixed-effects}
  \end{equation}
  where $\boldsymbol{\beta}$ are fixed-effect coefficients on the standardised feature vector
  $\mathbf{x}_{\text{plot}}$, $u_{\text{cpmt}}\sim\mathcal{N}(0,\sigma_u^2)$ is a compartment-level
  random intercept, and $\varepsilon_{\text{plot}}$ is residual error. This is a practical
  two-stage approximation to a full nonlinear mixed-effects refit of the CR curve itself, which
  was not implemented for this project. Fitted by restricted maximum likelihood (REML) for
  coefficient reporting; a separate variance-explained comparison between nested fixed-effects
  specifications uses maximum likelihood (ML) instead, since REML variance components are not
  directly comparable across models that differ in fixed-effects structure
  \citep{pinheiro2000}. % TODO: citation -- Pinheiro & Bates, no .bib entry yet.
  This model is fitted on the training partition only and has no held-out test evaluation step;
  its role is coefficient/variance reporting on the fitted training sample, not held-out
  prediction.
  % TODO: verify -- confirm the default random-effects structure used (random intercept only, no
  % random slope) against the fitting package's installed version, since no explicit random-effects
  % formula is set elsewhere in this description.
\end{itemize}

\subsection{Plot-specific growth-curve deviation (RQ3)}
\label{sec:rq3}

This avenue asks the same broad question at plot level: does a plot's \emph{own} fitted growth
ceiling deviate from what a static, external yield-class table would predict for it, and does
that deviation associate with environmental/management conditions. The deviation implemented here
is specifically in the fitted asymptote, with the curve's other shape parameters held fixed --
not a claim about the full growth trajectory.

\subsubsection{Target construction}
\label{sec:rq3-target}

For each plot, a per-plot asymptote $\hat y_{\max}^{\text{plot}}$ is fitted using \emph{only that
plot's own} repeated height/age observations -- never pooled across plots, unlike
Equation~\ref{eq:cr}. Because the curve's shape parameters are held fixed per plot (taken from
its own Forestry Commission yield-class lookup, not refitted), height is linear in the one
remaining unknown, $y_{\max}$, so the fit is closed-form weighted least-squares through the
origin rather than an iterative optimisation. The target is the difference between this
plot-specific fitted asymptote and the asymptote implied by the plot's own static yield-class
table entry,
\begin{equation}
\Delta y_{\max}^{\text{plot}} = \hat y_{\max}^{\text{plot}} - y_{\max}^{\text{yldc}}.
\label{eq:local-ymax-diff}
\end{equation}
As with RQ2, this target is built entirely from height/age survey data and a static external
lookup table, with no environmental variable involved in its construction. Each per-plot fit uses
only that plot's own small number of repeated observations (as few as the cohort minimum), so the
fitted asymptote carries more sampling uncertainty for plots with fewer surveys
(\S\ref{sec:limitations}).

\subsubsection{Global and spatially varying attribution models}
\label{sec:rq3-models}

Elastic Net and XGBoost are fitted as in \S\ref{sec:rq2-models}, on RQ3's own Set1--Set5. A third
model, GNNWR (geographically and neurally weighted regression), additionally allows coefficients
to vary continuously over space rather than assuming one global linear relationship,
\begin{equation}
\hat y_i = \beta_0(s_i) + \sum_j \beta_j(s_i)\,x_{ij},
\label{eq:gnnwr}
\end{equation}
where $s_i$ is plot $i$'s spatial location and each $\beta_j(s_i)$ is a locally varying
coefficient rather than a single global one \citep{du2020}. % TODO: citation -- Du et al. 2020, no .bib entry yet.
Concretely, a global OLS fit on the same predictors is taken first and frozen; a spatial weighting
network then takes, for each plot, the vector of Euclidean distances from that plot to every plot
in a reference set (by default the training partition), and outputs one multiplicative weight per
predictor plus an intercept weight. Each local coefficient $\beta_j(s_i)$ is that plot's weight
times the corresponding frozen global coefficient, so the network learns a spatially varying
re-weighting of one shared global regression rather than an unconstrained per-plot model; one
shared network produces every coefficient's weight jointly, not a separate network per coefficient.
There is no explicit distance-decay kernel or neighbourhood/bandwidth step -- weighting is learned
implicitly from the distance vector -- and held-out plots at evaluation are given their distance to
the training reference set, not a bespoke neighbour set of their own. The reference set is capped
at a fixed row limit for memory reasons, a disclosed deviation from using the full training
partition as every other model in this chapter does. This tests whether the attribution result
itself is spatially non-stationary, a question the global Elastic Net/XGBoost fits cannot answer.
Gain-based XGBoost importance and SHAP values are
computed after model fitting, over the full population rather than a held-out partition, purely to
interpret an already-chosen model; neither is used at any point to decide which variables enter a
feature set (\S\ref{sec:feature-screening}).
% TODO: interpreting SHAP on the same rows used to fit the model risks optimistic or unstable
% explanations even though SHAP does not feed back into feature selection -- consider out-of-fold
% or held-out SHAP, or state explicitly why full-population SHAP was judged acceptable here.

% FIGURE PLACEHOLDER (RQ3 target construction and attribution models, incl. GNNWR mechanism)
% Caption: "RQ3's asymptote-deviation target, fitted from each plot's own repeated observations
% against a static yield-class expectation, feeds three attribution models: global Elastic Net
% and XGBoost, and GNNWR, whose spatial weighting network re-weights a frozen global OLS fit
% using each plot's distances to a training reference set, producing plot-level local
% coefficients and predictions evaluated under the spatial-block split. RQ2's parallel but
% simpler pipeline (pooled residual -> Elastic Net/XGBoost/two-stage mixed-effects model) is
% described in prose and equations only (\S\ref{sec:rq2}), not duplicated in a figure here,
% since it has no comparable internal mechanism to unpack."
%
% Verified against the actual GNNWR implementation (third-party `gnnwr` package) before drawing:
% - input to the spatial weighting network is a vector of Euclidean distances from the focal
%   plot to every plot in the reference set (MinMax-scaled), NOT raw coordinates and NOT
%   coordinate differences;
% - no k-nearest-neighbour or distance-threshold step exists -- "neighbourhood" is implicit,
%   learned by the network from the full distance vector to the (capped) reference set;
% - the network does NOT output final coefficients directly: it outputs one multiplicative
%   weight per predictor (+ intercept), which is multiplied against a separately pre-fitted,
%   frozen global OLS coefficient to give the local coefficient beta_j(s_i);
% - one shared network with a multi-output head produces every coefficient's weight jointly, not
%   one network per coefficient;
% - held-out/test plots reuse the training reference set for their distance vector, not a
%   bespoke neighbour set built from their own location;
% - this is the standard GNNWR (Du et al. 2020) formulation, with one disclosed deviation: the
%   reference set is capped at a fixed row count (not the full training partition) for memory
%   reasons;
% - no geographical kernel or bandwidth parameter exists in the implementation -- do not draw one;
% - no rendered spatial coefficient map was found in the modelling pipeline, only a per-plot
%   coefficient table and a Moran's I residual-autocorrelation diagnostic -- label the coefficient
%   output as a "per-plot coefficient table", not a "coefficient map", unless a map-producing
%   script is confirmed separately.
%
% Mermaid sketch (for redrawing, not final layout):
% ```mermaid
% flowchart TD
%     A["Repeated heights, own plot only"] --> B["Closed-form per-plot<br/>asymptote fit"]
%     C["Yield-class asymptote<br/>(static lookup)"] --> D["RQ3 target: y_max diff"]
%     B --> D
%     F["Environmental/management<br/>predictors (Set1-5)"] --> G["Elastic Net"]
%     F --> H["XGBoost"]
%     F --> M["Global OLS fit<br/>(frozen coefficients)"]
%     D --> G
%     D --> H
%     D --> M
%     P["Plot location"] --> Q["Distance to every plot in<br/>reference set (train, capped)"]
%     Q --> R["Spatial weighting network<br/>(one shared trunk)"]
%     R --> S["Per-predictor weights w_j(s_i)"]
%     M --> T["Local coefficient<br/>beta_j(s_i) = w_j(s_i) x beta_j"]
%     S --> T
%     T --> U["Plot-level prediction"]
%     T --> V["Per-plot coefficient table"]
%     G --> W["Spatial-block holdout evaluation"]
%     H --> W
%     U --> W
% ```
%
% Content (prose fallback if not drawn from the mermaid sketch above): left-to-right, target
% construction (closed-form per-plot asymptote fit vs. static yield-class asymptote) feeding
% three model boxes -- Elastic Net, XGBoost, and an expanded GNNWR box showing the
% distance-to-reference-set input, spatial weighting network, per-predictor weights, and their
% multiplication against the frozen global OLS coefficients -- all three converging on a
% spatial-block holdout evaluation box. A small annotation: "association, not causation."
\begin{figure}[!htbp]
\centering
\fbox{\parbox{0.9\textwidth}{\centering\vspace{4cm}FIGURE PLACEHOLDER --- RQ3 target construction
and attribution models, see comment above source for exact content\vspace{4cm}}}
\caption{Plot-specific asymptote-deviation target construction (RQ3) and its three attribution
models, including the GNNWR spatial weighting mechanism.}
\label{fig:rq3-attribution}
\end{figure}

\section{Model evaluation, uncertainty and interpretation}
\label{sec:evaluation}

This section states what is measured and why, not what was found (Results chapter).

\subsection{Predictive and spatial diagnostics}
\label{sec:eval-diagnostics}

Every RQ1 model is scored on the same metric set, computed from residuals $e_i = H_i - \hat H_i$
(observed minus predicted height) over $N$ rows:
\begin{align*}
\text{MAE} &= \frac{1}{N}\sum_i |e_i|, &
\text{RMSE} &= \sqrt{\frac{1}{N}\sum_i e_i^2}, \\
R^2 &= 1 - \frac{\sum_i e_i^2}{\sum_i (H_i - \bar H)^2}, &
\text{Bias} &= \frac{1}{N}\sum_i e_i,
\end{align*}
where $\bar H$ is the mean observed height in that evaluation set; a positive bias indicates
systematic under-prediction. Metrics are also broken down by five age bands ($<30$, 30--40,
40--60, 60--80, $80+$ years) to check whether an aggregate score conceals age-specific weaknesses;
the maturity gate (plot-level age $\geq$ 30 at the 2023 reference survey) is applied at plot
level, so individual rows below 30 can still appear in a trajectory, and a poor score in the
$<30$ band reflects a small, non-representative sample rather than necessarily worse prediction.

Moran's I is computed on model residuals (not on raw features, and not as a variable-selection
criterion) to test whether spatial structure remains unexplained after fitting. This is a
diagnostic applied after a model is fitted, answering a different question from environmental
feature selection (\S\ref{sec:feature-sets}); it does not feed back into which variables enter a
model.

\subsection{Uncertainty and attribution stability}
\label{sec:eval-uncertainty}

Robustness of the RQ1 result to stochastic training variation is checked by a reseed of the
winning model across additional seeds, reported alongside the pooled-fold evaluation in the
Results chapter. The physics-weight ablation (\S\ref{sec:rq1-ablation}) plays the same role for
the physics/trajectory loss terms specifically: whether the guided network's advantage, if any,
depends on the weighting rather than the mechanism.

% TODO: state which of the following are actually implemented and where their settings live --
% bootstrap confidence intervals on RQ1 metrics; Elastic Net coefficient stability across folds/
% seeds; SHAP stability across folds/seeds for RQ3; spatial autocorrelation of the RQ2/RQ3
% attribution residuals themselves (as distinct from RQ1's residual diagnostic above). If any of
% these are not yet run, state that plainly rather than implying they are.

\section{Methodological assumptions and limitations}
\label{sec:limitations}

% Keep to 3-5 items. Draft points below, expand each into 2-4 sentences by hand.

\begin{itemize}
\setlength\itemsep{0pt}
\item \textbf{Spatial-block and buffer assumptions.} The 60\,m buffer is set from typical plot
  spacing, not from the measured spatial-autocorrelation range of the shared-curve residual
  (\S\ref{sec:splits}); compartment-level holdout is the primary leakage control, the buffer a
  secondary one.
\item \textbf{Environmental scale mismatch.} Terrain/wind/soil variables are raster-derived and
  assigned to point-scale plots; the raster resolution does not necessarily match the scale at
  which these conditions act on a plot.
  % TODO: confirm raster resolution(s) and the assignment method (nearest-cell, bilinear, zonal
  % mean) against the actual extraction code before stating this as fact.
\item \textbf{Feature-selection leakage or split dependence.} Ranking and VIF screening
  (\S\ref{sec:feature-screening}) determine Set1--Set5 before final evaluation; the open questions
  flagged there about which partition was used, and whether ranking was fold-specific, bear
  directly on how much confidence to place in the resulting sets.
\item \textbf{Uncertainty in plot-specific targets.} RQ3's per-plot asymptote
  (\S\ref{sec:rq3-target}) is fitted from as few as the cohort minimum of repeated observations
  per plot; plots with shorter survey histories contribute a noisier estimate of the same target
  as plots with longer histories, without that difference being explicitly weighted or flagged
  downstream.
\item \textbf{Association, not causation.} Both attribution avenues (\S\ref{sec:attribution})
  report statistical association between environmental/management conditions and departure from
  an expected growth curve; none of the designs support a causal claim.
\end{itemize}

\section{Reproducibility and implementation}
\label{sec:reproducibility}

All split assignments, network weight initialisation, and Elastic Net/XGBoost cross-validation
folds use a fixed seed (42) unless explicitly swept as part of a stated robustness check (e.g. the
RQ1 winner-model reseed check, \S\ref{sec:eval-uncertainty}). Every model reads its feature list
from a single generated manifest, so a reported set's exact column membership can always be
checked against that source rather than a hand-copied list.

Detailed hyperparameter grids and selected values (Appendix~\ref{app:hyperparameters}), full
Set1--Set5 column lists and per-variable screening diagnostics (Appendix~\ref{app:feature-sets}),
and secondary sensitivity settings not reported in the main text -- dropout-rate/architecture-width
sweeps, the narrow-gap temporal split's own results (Appendix~\ref{app:sensitivity}) -- are kept
out of this chapter and pointed to at the relevant section above rather than repeated here.

% TODO: verify -- confirm whether RQ2/RQ3's own model-fitting scripts read the same feature-set
% manifest as RQ1, or still require a separate wiring step, before this chapter's Set1-5
% description is presented as equally accurate for all three research questions.
